\section{Results and Discussion}
\label{sec:results}

\begin{figure*}[t]
\centering
\includegraphics[width=0.96\textwidth]{parameter_history.png}
\caption{Accepted-state objective and parameter history for the preliminary five-parameter run. The orange ring marks the lowest-loss saved state; all 16 retained states are shown without smoothing.}
\label{fig:parameters}
\end{figure*}

The stored Episode~18 diagnostics use distinct fixed configurations, so their losses and parameter values are not pooled across runs. Figure~\ref{fig:parameters} reports the preliminary five-parameter history, whereas Fig.~\ref{fig:matched} shows outputs from a separate strictly evaluated run.

\subsection{Preliminary Five-Parameter Fitting}

The saved history contains the initial state and 15 accepted updates. The dimensionless partial-view training objective decreased from $0.42455$ to $0.16229$, a reduction of $61.8\%$. Most of the reduction occurred by update 7; updates 10--15 changed the objective from $0.16322$ to $0.16229$, a further $0.57\%$, indicating a late-stage plateau in the saved history. The lowest-loss saved state is
\begin{equation}
\begin{split}
E&=8.80\ \mathrm{kPa},\quad \nu=0.487,\quad
\eta=30.4\ \mathrm{Pa\,s},\\
\sigma_{\min}&=0.8626,\qquad \sigma_{\max}=1.0671.
\end{split}
\label{eq:preliminary_parameters}
\end{equation}
The variables move together during optimization: viscosity increases, both stretch limits shift downward with little change in interval width, and $\nu$ approaches the upper part of its allowed range. This coordinated change is consistent with compensation among constitutive effects under one trajectory. The values in~\eqref{eq:preliminary_parameters} are therefore a preliminary effective-parameter result, not final material calibration.

\subsection{Visible-Geometry Evaluation}

A separate five-parameter run with 20 accepted updates supplies Fig.~\ref{fig:matched}. Its selected state reached the $E=300\,$kPa upper bound and the $\eta=0$ lower bound. A strict full-resolution evaluator equally averages normalized depth-change error, one minus mask intersection-over-union, directed observed-to-simulation distance, and a missing-coverage penalty. Its mean score, for which lower is better, was $0.2821$ compared with $0.3416$ for frozen initial geometry on frames 1--60. On frames 61--97 of the same replay, the corresponding scores were $0.3939$ and $0.8265$; all 60 fitting frames and 37 continuation frames were valid. The run therefore improves visible-geometry prediction relative to a static-object baseline during the recorded motion, but the parameter bounds prevent interpreting it as calibrated physical recovery.

Figure~\ref{fig:matched} shows qualitative agreement in gross visible geometry through $1.00\,$s. By $2.00\,$s, the simulated visible geometry differs markedly from the recording, and little subsequent deformation is apparent in the simulation. This comparison is descriptive because the rows show different modalities and simulated depth does not include tool occlusion.

\subsection{Limitations}

A single depth view leaves hidden volume and occluded contact unmeasured, and the evaluator does not render the spatulas into simulated depth. Both reported runs use the fixed table-aligned calibration described in Sec.~\ref{sec:method}; the estimates nevertheless remain conditional on residual camera, floor, and tool-registration uncertainty, reconstructed volume and density, contact law and friction, numerical resolution, and loss smoothing. Near-incompressible $\nu$, elasticity, viscosity, and stretch clipping can also compensate for one another under a single trajectory.

No tool force or torque measurements were recorded. The optimizer therefore distinguishes contact and constitutive effects only through visible motion and geometry, although different combinations can produce similar depth trajectories while implying different contact forces. With one sequence and no comparison against competing systems on common data, the experiment does not establish state-of-the-art identification accuracy or uncertainty across repeated preparations. It demonstrates end-to-end fitting and measurable improvement over frozen geometry in the same recorded replay.

\FloatBarrier
